% !TEX root = ../main.tex

%%% ============================================================
\section{Kết luận}
%%% ============================================================

\subsection{Kết luận \& Hướng phát triển}
\begin{frame}{Kết luận, hạn chế và hướng phát triển}
  \framesubtitle{Đã làm được gì, và chưa làm được gì?}
  \begin{columns}[t]
    \begin{column}{0.48\textwidth}
      \textbf{Kết quả đạt được}
      \begin{itemize}[<+->]\ilist
        \item \kw{146 skill}, \kw{57 operator} $\rightarrow$ \kw{1760 trajectory} / \kw{9007 mẫu SFT}, 0 nhãn người.
        \item \kw{13,3\% $\rightarrow$ 24,6\%} ($p<0{,}001$), ngang SOTA ${<}10$B.
        \item \kw{0,0308 \$/trajectory} --- rẻ hơn 3,3 lần.
        \item Bàn giao mã nguồn, dataset, checkpoint.
      \end{itemize}
    \end{column}
    \hfill
    \begin{column}{0.48\textwidth}
      \textbf{Hạn chế}
      \begin{itemize}[<+->]\isub
        \item Điểm nghẽn ở tầng thực thi (timeout DOM).
        \item Chưa chạy ablation so với khám phá phẳng.
        \item \textit{Phạm vi}: chuyên biệt theo website; cross-app ngoài phạm vi (PL1).
      \end{itemize}

      \gap
      \uncover<8->{
        \textbf{Hướng phát triển}
        \begin{itemize}\isub
          \item Tự phục hồi sau lỗi (\textit{error recovery}).
          \item Khai thác tín hiệu thất bại bằng RL / DPO.
          \item Chưng cất xuống SLM 1--3B.
        \end{itemize}
      }
    \end{column}
  \end{columns}
\end{frame}

\begin{frame}[plain]
  \centering
  \vfill
  {\usebeamerfont{title}\usebeamercolor[fg]{title}\Large\textbf{XIN CẢM ƠN HỘI ĐỒNG}}\\[0.3cm]
  {\normalsize Nhóm nghiên cứu sẵn sàng lắng nghe câu hỏi và góp ý}
  \vfill
  \footnotesize
  \begin{tabular}{ll}
    \textbf{Sinh viên thực hiện:} & Trương Bảo Kiệt (22127223) - Hồ Tiến Phát (22127321) \\
    \textbf{GVHD:} & PGS.TS. Nguyễn Văn Vũ \\
    \textbf{Đơn vị:} & Bộ môn CNPM, Khoa CNTT - Trường ĐH KHTN, ĐHQG-HCM
  \end{tabular}
  \vfill
\end{frame}


